% note.tex -- Smallest enclosing balls by an active-set method on the dual.
% Build:  pdflatex note.tex   (twice)   or   tectonic note.tex

\documentclass[10pt,a4paper]{article}

\usepackage[T1]{fontenc}
\usepackage{amsmath,amssymb,amsthm}
\usepackage{mathtools}
\usepackage{booktabs}
\usepackage{xcolor}
\usepackage[margin=2.4cm]{geometry}
\usepackage[colorlinks=true,allcolors=blue!50!black]{hyperref}

\theoremstyle{plain}
\newtheorem{proposition}{Proposition}
\newtheorem{lemma}[proposition]{Lemma}
\theoremstyle{remark}
\newtheorem*{remark}{Remark}

\DeclareMathOperator{\conv}{conv}
\DeclareMathOperator{\aff}{aff}
\DeclareMathOperator*{\argmax}{arg\,max}
\newcommand{\R}{\mathbb{R}}
\newcommand{\norm}[1]{\lVert #1 \rVert}
\newcommand{\T}{^{\mathsf{T}}}

\setlength{\parskip}{0pt}
\title{\vspace{-1.2cm}\bfseries Smallest enclosing balls by an active-set method on the dual}
\author{Thomas Schmelzer}
\date{\today}

\begin{document}
\maketitle
\vspace{-0.6cm}

\begin{abstract}
\noindent
We record the dual quadratic program of the smallest enclosing ball problem, the
active-set method that solves it, and the invariance properties an
implementation must preserve. The method terminates at an exact vertex of the
dual feasible set rather than at an interior-point tolerance, costs one
$k\times k$ dense solve per iteration with $k\le d$, and returns the dual
weights, which are the conic dual of the second-order cone formulation after one
rescaling and hence a certificate the caller can check independently.
\end{abstract}

\section{The problem and its dual}

For $S=\{p_1,\dots,p_n\}\subset\R^d$ the smallest enclosing ball is the solution of
\begin{equation}
  \label{eq:primal}
  \min_{x,r}\ r \quad\text{s.t.}\quad \norm{p_i-x}\le r,\ i=1,\dots,n,
\end{equation}
a second-order cone program in $(r,x)\in\R^{1+d}$ with $n$ blocks
$(r,\,p_i-x)\in\mathcal{Q}^{d+1}$. Existence follows from coercivity of
$f(x)=\max_i\norm{p_i-x}$ and uniqueness from strict convexity of $f^2$ along any
segment; we write $(R,x^\star)$ for the optimum.

Let $\Delta_n=\{u\in\R^n: u\ge0,\ \mathbf{1}\T u=1\}$, and for $u\in\Delta_n$ put
$x(u)=\sum_i u_ip_i$ and
\begin{equation}
  \label{eq:phi}
  \varphi(u)=\sum_i u_i\norm{p_i}^2-\norm{x(u)}^2 ,
\end{equation}
a concave quadratic. Everything below rests on one identity, valid for all
$y\in\R^d$ and $u\in\Delta_n$:
\begin{equation}
  \label{eq:id}
  \sum_i u_i\norm{p_i-y}^2=\varphi(u)+\norm{x(u)-y}^2 .
\end{equation}

\begin{proposition}
  \label{prop:duality}
  $R^2=\max_{u\in\Delta_n}\varphi(u)$, and $x^\star=x(w)$ for any maximiser $w$.
\end{proposition}

\begin{proof}
  If $B(y,r)\supseteq S$ then the left side of~\eqref{eq:id} is at most $r^2$, so
  $\varphi(u)\le r^2$ for every $u$; taking the optimum gives
  $\max\varphi\le R^2$. Conversely let $T=\{p_i:\norm{p_i-x^\star}=R\}$. If
  $x^\star\notin\conv T$, separate: there is $v$ with $v\T(p-x^\star)>0$ on $T$,
  and $\norm{p-(x^\star+tv)}^2=R^2-2tv\T(p-x^\star)+t^2\norm{v}^2<R^2$ for small
  $t>0$, while the finitely many points off $T$ stay strictly inside --- so the
  radius decreases, a contradiction. Hence $x^\star=\sum_i w_ip_i$ for some
  $w\in\Delta_n$ supported on $T$, and $y=x^\star$ in~\eqref{eq:id} gives
  $\varphi(w)=\sum_iw_i\norm{p_i-x^\star}^2=R^2$. The same substitution forces
  $\norm{x(w)-x^\star}^2\le0$ for any maximiser.
\end{proof}

Since $\partial\varphi/\partial u_i=\norm{p_i-x(u)}^2-\norm{x(u)}^2$, the
first-order conditions for $\max_{\Delta_n}\varphi$ read, after cancelling the
common term,
\begin{equation}
  \label{eq:kkt}
  \norm{p_i-x}=R\ \ (w_i>0),\qquad \norm{p_i-x}\le R\ \ (w_i=0),
\end{equation}
the geometric certificate: the points carrying weight lie on the boundary sphere
and the rest are enclosed. By Carath\'eodory the support may be taken of size at
most $d+1$, so~\eqref{eq:primal} is determined by at most $d+1$ of the $n$ points
and the problem is combinatorial in the choice of that subset.

\section{The method}

Maintain a working set $F$ and weights $u>0$ on it summing to one.

\paragraph{Subproblem.} Maximising $\varphi$ over weights supported on $F$ with no
sign restriction is, by~\eqref{eq:kkt} without the inequalities, the requirement
that $x$ be equidistant from $\{q_j\}_{j\in F}$ and lie in their affine hull.
With $D$ the $(m-1)\times d$ matrix of edges $e_j=q_j-q_0$ and $x=q_0+D\T y$,
\begin{equation}
  \label{eq:sub}
  (DD\T)\,y=\tfrac12\bigl(\norm{e_1}^2,\dots,\norm{e_{m-1}}^2\bigr)\T,
  \qquad w=\Bigl(1-\textstyle\sum_j y_j,\ y\Bigr),
\end{equation}
a dense system of order $m-1\le d$, non-singular exactly when $F$ is affinely
independent. This is the circumcentre of the face, expressed in barycentric
coordinates.

\paragraph{Dropping.} If $\min_j w_j<0$ the circumcentre lies outside
$\conv\{q_j\}$ and the subproblem solution is dual-infeasible. Move along
$s=w-u$ by the ratio test
\begin{equation}
  \label{eq:ratio}
  \alpha=\min\{-u_i/s_i:\ s_i<0\},
\end{equation}
which drives at least one weight to zero; remove it from $F$.

\paragraph{Degeneracy.} If $F$ is affinely dependent, $DD\T$ is singular. The
directions $s$ with $\mathbf{1}\T s=0$ and $\sum_i s_iq_i=0$ --- equivalently
$D\T s_{1:}=0$ after eliminating $s_0$ --- leave $x(u)$ fixed, so
$\norm{x(u)}^2$ is frozen and $\varphi(u+ts)=\varphi(u)+t\sum_is_i\norm{p_i}^2$
is linear: the subproblem is unbounded. Project the gradient onto that null space
and apply~\eqref{eq:ratio}, shrinking $F$ by one.

\paragraph{Adding.} Otherwise accept $u\leftarrow w$, set $x=x(u)$ and
$R=\max_{i\in F}\norm{p_i-x}$, and let $k=\argmax_i\norm{p_i-x}$. If
$\norm{p_k-x}\le R$ then~\eqref{eq:kkt} holds and $(R,x,u)$ is optimal; otherwise
append $k$ to $F$ with weight zero. By~\eqref{eq:gradient-note} below this is
also the steepest-ascent choice among the fixed variables.

\begin{center}
\fbox{\begin{minipage}{0.95\textwidth}
\small
\textbf{Input} $p_1,\dots,p_n$. Recentre on $\bar p$. Set $F=\{k\}$, $k$ farthest
from $p_1$; $u_k=1$.\\
\textbf{Repeat:} if $F$ affinely dependent, take the null-space descent step and
drop; else solve~\eqref{eq:sub}. If $\min w<0$, take the ratio step~\eqref{eq:ratio}
and drop. Else set $u=w$, $x=x(u)$, $R=\max_{i\in F}\norm{p_i-x}$; if no point
exceeds $R$, return $(R,\,x+\bar p,\,u)$; else drop zero weights and append
$\argmax_i\norm{p_i-x}$ with weight $0$.
\end{minipage}}
\end{center}

Each iteration either strictly increases $\varphi$ or strictly shrinks $F$, and
there are finitely many working sets, so the method is finite; as with any
active-set method a degenerate zero-length step requires an anti-cycling rule in
exact arithmetic, and in practice an iteration cap serves as a safety net. Note
that
\begin{equation}
  \label{eq:gradient-note}
  \partial\varphi/\partial u_k-\partial\varphi/\partial u_i
    =\norm{p_k-x}^2-\norm{p_i-x}^2,
\end{equation}
so the farthest violator is the fixed variable with the largest reduced gradient.

\section{Invariance}

Two properties are load-bearing and easily lost. Both are consequences of
refusing any tolerance tied to coordinates of magnitude one.

\emph{Origin invariance.} Recentre on $\bar p$ before iterating, and form
$\norm{p-x}^2$ as $\sum_k(p_k-x_k)^2$ rather than
$\norm{p}^2-2p\T x+\norm{x}^2$. The expanded form cancels quantities of size
$\norm{x}^2$ to produce an answer of size $R^2$, losing every digit by which
$\norm{x}$ exceeds $R$; for a cloud of unit extent at distance $10^6$ the
rounding floor of a squared distance is then of the order of the radius itself.

\emph{Scale invariance.} Test affine dependence on the rank of $D$, not on
$[\mathbf{1};Q\T]$: the row of ones has unit entries, dominates the singular
values, and would declare any cloud of extent below $\varepsilon$ degenerate.
Size the feasibility slack of the stopping test as a relative tolerance plus an
absolute floor proportional to $\varepsilon\max_i\norm{p_i}^2$, and normalise the
null-space direction before comparing anything derived from it --- the gradient
carries units of length$^2$, the weight tolerance does not.

\section{Cost, and the dual as a certificate}

Per iteration the method performs one $O(nd)$ sweep of squared distances and one
dense solve of order at most $d$; nothing of size $n$ is ever factorised. This is
the structural difference from an interior-point method applied
to~\eqref{eq:primal}, which factors an $n(d+1)$-row system at every step. The
method also identifies the support combinatorially and solves that face exactly,
rather than approaching it from the interior and leaving the caller to decide
which slacks of size $10^{-9}$ were meant to be zero.

The weights are the conic dual of~\eqref{eq:primal} after one rescaling: the
block for point $i$ is
\begin{equation}
  \label{eq:conic-dual}
  z_i=u_i\Bigl(1,\ -\frac{p_i-x}{R}\Bigr)\in\mathcal{Q}^{d+1},
\end{equation}
whose tail has norm $u_i\norm{p_i-x}/R$, equal to its head exactly on the
support, and which pairs with the slack $s_i=(R,\,p_i-x)$ as
$s_i\T z_i=(u_i/R)(R^2-\norm{p_i-x}^2)=0$ --- complementarity, each factor
vanishing where the other does not. So a solver built on this method returns a
standard-form primal--dual pair, and a caller can verify optimality
from~\eqref{eq:kkt} at the cost of one pass over the data, without re-solving.

\begin{remark}[Substitutability]
  This is what makes the method usable behind a conic solver's interface rather
  than only as a special-purpose routine. The accompanying implementation
  dispatches on the cone list: zero and non-negative blocks go to a dense dual
  active-set QP method~\cite{goldfarb1983}, and equally sized second-order blocks
  matching the shape of~\eqref{eq:primal} go to the method above. A program
  carrying the right cones but a different matrix is refused rather than
  answered, which requires reconstructing the point cloud from $(A,b)$ and
  rebuilding the program for comparison. An indefinite objective is likewise
  refused rather than regularised into definiteness, since $P+\epsilon I$ answers
  a different question.
\end{remark}

\section{A numerical comparison}

Table~\ref{tab:bench} compares the method with Clarabel~0.11.1 applied
to~\eqref{eq:primal} on Gaussian clouds, measuring the relative enclosure error
$\max_i\norm{p_i-x}/R-1$ --- positive values mean the returned ball fails to
contain the cloud. Both agree on $R$ to the digits shown. The active-set figures
are $0$ or a single ulp, since the returned $R$ is by construction the exact
maximum distance over the support of a centre computed from~\eqref{eq:sub};
Clarabel's are its interior-point residuals, four to five orders larger. The
timing gap grows with $n$ because only the $O(nd)$ sweep depends on it.

\begin{table}[h]
\centering
\small
\begin{tabular}{rr rr rr}
  \toprule
  & & \multicolumn{2}{c}{enclosure error} & \multicolumn{2}{c}{time (s)} \\
  \cmidrule(lr){3-4}\cmidrule(lr){5-6}
  $n$ & $d$ & active set & Clarabel & active set & Clarabel \\
  \midrule
  $10^3$ &  2 & $0$                  & $5.3\cdot10^{-12}$ & 0.001 & 0.009 \\
  $10^3$ & 10 & $0$                  & $9.4\cdot10^{-12}$ & 0.001 & 0.056 \\
  $10^4$ &  3 & $2.2\cdot10^{-16}$   & $2.0\cdot10^{-12}$ & 0.001 & 0.141 \\
  $10^4$ & 20 & $-1.1\cdot10^{-16}$  & $3.5\cdot10^{-14}$ & 0.006 & 0.948 \\
  $10^5$ &  3 & $-2.2\cdot10^{-16}$  & $1.7\cdot10^{-11}$ & 0.008 & 2.225 \\
  $10^5$ & 10 & $0$                  & $-7.2\cdot10^{-12}$ & 0.014 & 6.595 \\
  \bottomrule
\end{tabular}
\caption{Standard normal clouds, one seed per row; Apple M4 Pro, CPython 3.12.13,
NumPy 2.5.1, Clarabel 0.11.1 at its default tolerances. Times are a single run
and are indicative only; the enclosure errors are not, being deterministic
properties of the returned pair.}
\label{tab:bench}
\end{table}

The comparison is not entirely fair to Clarabel, which is built for large sparse
programs of arbitrary shape and is here being asked for a dense problem of one
particular shape whose combinatorial structure the other method exploits. That
is the point: the specialised method is worth keeping precisely where the
structure is known.

\section{Relation to the minimum-norm point problem}

If all points have equal norm $\rho$ --- for instance after the standard lift of
a kernel problem onto a sphere --- then~\eqref{eq:phi} reduces to
$\varphi(u)=\rho^2-\norm{x(u)}^2$, so maximising $\varphi$ over $\Delta_n$ is
minimising $\norm{x}$ over $\conv S$: the minimum-norm point problem solved by
Wolfe's algorithm and by the Frank--Wolfe family. The working set of
Section~2 is then Wolfe's corral, the subproblem~\eqref{eq:sub} is his affine
minimiser, and the drop step is his major cycle. The general case is the same
algorithm with the linear term $\sum_iu_i\norm{p_i}^2$ restored, which is what
makes the subproblem a circumcentre rather than a projection.

\section{Relation to Welzl's algorithm}

Welzl's randomised incremental method~\cite{welzl1991} computes the same object
and, like the method above, terminates at an exact basis rather than at a
tolerance. The two differ in how that basis is searched.

Welzl's search is combinatorial: it recurses over subsets, drawing a point $p$ at
random, solving without it, and re-solving with $p$ forced onto the boundary when
$p$ falls outside the ball just computed. The permutation drives the recursion,
the only arithmetic is the primitive that circumscribes at most $d+1$ points, and
the analysis is that of LP-type problems. The method of Section~2 searches
instead by ascent on $\varphi$: the weights say what to do next, a negative
barycentric coordinate calling for a drop and the farthest violator for an
addition, so the search is monotone and local, with add \emph{and} drop pivots,
and it is deterministic.

Three consequences are worth stating, the first of them not in this method's
favour. Welzl's algorithm runs in expected $O(n)$ time for fixed $d$, and nothing
comparable is claimed here --- finite termination, with the usual caveat about
degenerate steps, is all Section~2 establishes. But that bound holds $d$ fixed,
and its constant grows like $d!$, which is what confines the method in practice
to low dimension; the cost per iteration here is $O(nd)$ plus a solve of order at
most $d$, which is why the $d=20$ row of Table~\ref{tab:bench} is unremarkable.
This is the trade that motivated the support-set framework of Fischer, G\"artner
and Kutz~\cite{fischer2003}, to which the present method belongs.

Second, the shape of the work differs as much as its amount. The inner operation
here is one sweep of squared distances over all points --- a single vectorised
kernel, and the only place $n$ appears --- while Welzl's is a membership
predicate evaluated per recursion step, at a recursion depth of $O(n)$. On the
array-oriented runtimes where implementations of this problem now mostly live,
that difference dominates the constant factors on both sides.

Third, the two return different certificates. Welzl's recursion yields the ball
and its basis; the method above yields the barycentric weights as well, which
are what make~\eqref{eq:kkt} checkable in one pass and what
become~\eqref{eq:conic-dual} when a conic interface asks for a dual. The weights
can of course be recovered from a basis by solving~\eqref{eq:sub}, but they are
not a by-product of the search there, whereas here they \emph{are} the search.
The same asymmetry makes a support set a natural warm start for a perturbed
cloud, which a recursion over a fresh random permutation has no way to accept.

G\"artner's robust implementation~\cite{gaertner1999} addresses the numerical
fragility of the recursion's predicates, and should be read alongside the
original; the invariance requirements of Section~3 are this method's answer to
the same question.

\section{Notes}

The problem is Sylvester's~\cite{sylvester1857}; Elzinga and
Hearn~\cite{elzinga1972} gave an early finite algorithm of this type, and
Section~7 places it against the randomised incremental line of work. The
support-set framework of~\cite{fischer2003} also covers the smallest ball
enclosing a set of \emph{balls}; the method above is stated and implemented here
for points only. The QP machinery is standard~\cite{nocedal2006}. An
implementation of the method, of the Goldfarb--Idnani core, and of the conic
front end is available at \url{https://github.com/tschm/cvxball}.

\begin{thebibliography}{9}
\small
\bibitem{elzinga1972} J.~Elzinga and D.~W.~Hearn. Geometrical solutions for some
  minimax location problems. \emph{Transportation Science} 6(4):379--394, 1972.
\bibitem{fischer2003} K.~Fischer, B.~G\"artner and M.~Kutz. Fast
  smallest-enclosing-ball computation in high dimensions. \emph{Proc.\ ESA},
  630--641, 2003.
\bibitem{gaertner1999} B.~G\"artner. Fast and robust smallest enclosing balls.
  \emph{Proc.\ ESA}, 325--338, 1999.
\bibitem{goldfarb1983} D.~Goldfarb and A.~Idnani. A numerically stable dual
  method for solving strictly convex quadratic programs. \emph{Mathematical
  Programming} 27:1--33, 1983.
\bibitem{nocedal2006} J.~Nocedal and S.~J.~Wright. \emph{Numerical
  Optimization}, 2nd ed. Springer, 2006.
\bibitem{sylvester1857} J.~J.~Sylvester. A question in the geometry of
  situation. \emph{Quarterly Journal of Pure and Applied Mathematics} 1:79, 1857.
\bibitem{frank1956} M.~Frank and P.~Wolfe. An algorithm for quadratic
  programming. \emph{Naval Research Logistics Quarterly} 3:95--110, 1956.
\bibitem{welzl1991} E.~Welzl. Smallest enclosing disks (balls and ellipsoids).
  \emph{New Results and New Trends in Computer Science}, LNCS 555, 359--370, 1991.
\bibitem{wolfe1976} P.~Wolfe. Finding the nearest point in a polytope.
  \emph{Mathematical Programming} 11:128--149, 1976.
\end{thebibliography}

\end{document}
